\chapter{Basic Airway Management}

\section*{What Is This Test or Procedure?}
Basic airway management is the use of simple maneuvers and devices to open and maintain a patent airway in an unconscious or obstructed patient. It includes head-tilt and jaw-thrust positioning, oropharyngeal and nasopharyngeal airway insertion, and bag-valve-mask ventilation.
\section*{Alternative Names}
Airway opening maneuvers; basic life support airway; BVM ventilation; oropharyngeal airway insertion
\section*{Principle}
In the unconscious patient the tongue and soft tissues fall backward and obstruct the pharynx. Head-tilt and chin-lift or jaw-thrust lift the tongue forward, and oropharyngeal or nasopharyngeal airways hold the pharynx open. Bag-valve-mask ventilation then delivers breaths through the maintained airway.
\section*{History}
Airway management evolved from the anesthesia techniques of the late 19th century, and the bag-valve-mask device and standardized airway adjuncts became central to resuscitation and emergency care in the mid-20th century.
\section*{Why This Test or Procedure Is Ordered}
Performed in any patient with an obstructed airway, apnea, or impaired consciousness, during cardiopulmonary resuscitation, and as the first step before advanced airway placement such as endotracheal intubation.
\section*{Is This a Primary or Secondary Test or Procedure}
Primary lifesaving procedure for airway obstruction and a bridge to advanced airway management.
\section*{How This Test or Procedure Is Performed}
The patient is positioned supine, and the head-tilt/chin-lift (or jaw-thrust if cervical injury is suspected) is applied to open the airway. Secretions and foreign bodies are cleared, an appropriately sized airway adjunct is inserted, and ventilation is delivered with a bag-valve-mask at an appropriate rate and volume, checking for chest rise.
\section*{What Preparation Is Needed}
No specific preparation beyond suction, airway adjuncts of several sizes, a bag-valve-mask with reservoir, and supplementary oxygen. Cervical spine protection is maintained when trauma is suspected.
\section*{Contraindications and Precautions}
Oropharyngeal airways are avoided in patients with an intact gag reflex, and nasopharyngeal airways are avoided with facial fractures or suspected basilar skull fracture. Precautions include correct sizing and reassessment of effectiveness.
\section*{Risks}
Gastric insufflation with resultant aspiration, inadequate ventilation, oropharyngeal injury, and nasal bleeding from nasopharyngeal insertion.
\section*{What Happens After the Test or Procedure}
The patient is continuously monitored with pulse oximetry, capnography, and clinical assessment, and ventilation is continued until spontaneous breathing or definitive airway management is established.
\section*{Understanding Your Results}
Effective airway management is confirmed by chest rise, adequate oxygen saturation, and end-tidal carbon dioxide on capnography. Failure to ventilate prompts repositioning, adjunct adjustment, or advanced airway placement.
\section*{Normal Results}
A patent airway with visible chest expansion, audible breath sounds, and stable or improving oxygenation.
\section*{Follow-up Tests and Procedures}
Patients who remain unconscious or fail basic airway management proceed to endotracheal intubation, and those who recover are observed for aspiration, airway injury, and the underlying cause.
\section*{References}
\begin{enumerate}
  \item MedlinePlus. Basic life support and CPR. U.S. National Library of Medicine; 2025.
  \item Current Medical Diagnosis and Treatment. McGraw-Hill; 2025.
\end{enumerate}

\clearpage
